\begin{enumerate}
	\item $v = \dfrac{d}{t}=\dfrac{950}{10}= 95$.

La vitesse moyenne du trajet est de 95~km/h.

	\item Les panneaux possédant un centre de symétrie sont les panneaux 1, 5 et 6.

	\begin{center}
		\begin{tblr}{colspec={*{6}{X[c]}}, hlines, vlines}
			Panneau 1 & Panneau 2 & Panneau 3 & Panneau 4 & Panneau 5 & Panneau 6\\
			\includegraphics[width=1.5cm]{sens_interdit.eps} &
			\includegraphics[width=1.5cm]{interdit_dangereux.eps} &
			\includegraphics[width=1.5cm]{interdit_doubler.eps} &
			\includegraphics[width=1.5cm]{priorite_droite.eps} &
			\includegraphics[width=1.5cm]{interdit_arret.eps} &
			\includegraphics[width=1.5cm]{route_prioritaire.eps}\\
			oui & non & non & non & oui & oui\\
		\end{tblr}
	\end{center}

	\item
	\begin{enumerate}
		\item $2\times 950 = \np{1900}$.
		
La distance totale parcourue est de \np{1900}~km.

		\item Carburant consommé : $\dfrac{\np{1900}}{100}\times 5,5 = 19\times 5,5 = 104,5$~L.

Coût du carburant : $104,5 \times 1,70 = 177,65$~\euro.

Coût des péages : $65 + 65 = 130$~\euro.

Coût total : $177,65 + 130 = 307,65$~\euro{}
		
et $307,65 < 311$.

Le budget de 311~\euro{} est suffisant.
	\end{enumerate}
\end{enumerate}
